% Test: la anotación al margen de \guia dentro de un bloque display
% (\[...\]) debe caer en la MISMA página que el bloque que anota, incluso
% cuando ese bloque queda justo en el borde de una página. Antes de la
% corrección, \@guiamargnote encolaba la anotación y la vaciaba recién al
% cerrar el "\]" (\@guia@flushqueue); si ese cierre caía justo donde el
% "page builder" ya había cortado la página, la anotación terminaba
% flotando en la página SIGUIENTE a la del bloque display (ver README,
% sección "Anotación de \guia en la página siguiente al bloque display").
% El relleno de abajo (texto repetido sin significado) empuja el bloque
% "Diferenciando..." justo al borde de la página 1/2, verificado con esta
% clase y motor; si el bug reaparece, "Diferenciando" y los "+1" quedarán
% en páginas distintas.
\PassOptionsToClass{guia}{emate-ucr}
\documentclass{emate-ucr}
\curso{Test}
\encabezado{Test guia pagebreak}

\begin{document}
\imprimirtitulo

\begin{ejercicio}[1]
  Enunciado de relleno.
  \begin{solucion}
    \noindent
    Texto de relleno para llenar la pagina y forzar un corte justo antes del bloque display. Texto de relleno para llenar la pagina y forzar un corte justo antes del bloque display. Texto de relleno para llenar la pagina y forzar un corte justo antes del bloque display. Texto de relleno para llenar la pagina y forzar un corte justo antes del bloque display. Texto de relleno para llenar la pagina y forzar un corte justo antes del bloque display. Texto de relleno para llenar la pagina y forzar un corte justo antes del bloque display. Texto de relleno para llenar la pagina y forzar un corte justo antes del bloque display. Texto de relleno para llenar la pagina y forzar un corte justo antes del bloque display. Texto de relleno para llenar la pagina y forzar un corte justo antes del bloque display. Texto de relleno para llenar la pagina y forzar un corte justo antes del bloque display. Texto de relleno para llenar la pagina y forzar un corte justo antes del bloque display. Texto de relleno para llenar la pagina y forzar un corte justo antes del bloque display. Texto de relleno para llenar la pagina y forzar un corte justo antes del bloque display. Texto de relleno para llenar la pagina y forzar un corte justo antes del bloque display. Texto de relleno para llenar la pagina y forzar un corte justo antes del bloque display. Texto de relleno para llenar la pagina y forzar un corte justo antes del bloque display. Texto de relleno para llenar la pagina y forzar un corte justo antes del bloque display. Texto de relleno para llenar la pagina y forzar un corte justo antes del bloque display. Texto de relleno para llenar la pagina y forzar un corte justo antes del bloque display. Texto de relleno para llenar la pagina y forzar un corte justo antes del bloque display. Texto de relleno para llenar la pagina y forzar un corte justo antes del bloque display. Texto de relleno para llenar la pagina y forzar un corte justo antes del bloque display. Texto de relleno para llenar la pagina y forzar un corte justo antes del bloque display. Texto de relleno para llenar la pagina y forzar un corte justo antes del bloque display. Texto de relleno para llenar la pagina y forzar un corte justo antes del bloque display. Texto de relleno para llenar la pagina y forzar un corte justo antes del bloque display. Texto de relleno para llenar la pagina y forzar un corte justo antes del bloque display. Texto de relleno para llenar la pagina y forzar un corte justo antes del bloque display. Texto de relleno para llenar la pagina y forzar un corte justo antes del bloque display. Texto

    Diferenciando con respecto a $p_x$:
    \[
      \guia[1]{\frac{1}{z}\,\frac{\partial z}{\partial p_x} = \frac{3}{p_x}}
      \guia[1][-3\baselineskip]{\implies
      \frac{\partial z}{\partial p_x} = \frac{3z}{p_x}}.
    \]
  \end{solucion}
\end{ejercicio}

\end{document}

%% Local Variables:
%% mode: LaTeX
%% TeX-master: t
%% End:
